\documentclass[a4paper]{article}

\usepackage[spanish]{babel}
\usepackage{graphicx}
\usepackage{amsmath, amssymb}
\usepackage[margin=2cm]{geometry}
\usepackage{fancyhdr}
\usepackage{enumerate}
\usepackage[shortlabels]{enumitem}
\usepackage{parskip}
\usepackage[most]{tcolorbox}
\usepackage[hidelinks]{hyperref}
\usepackage{float}
\usepackage{booktabs}



% cabecera
\pagestyle{fancy}
\fancyhead[l]{Fisica Matemática I}
\fancyhead[c]{Notas de Clase Parte 1}
\fancyhead[r]{\today}
\fancyfoot[c]{\thepage}
\renewcommand{\headrulewidth}{0.2pt} % linea horizontal

\begin{document}

\begin{titlepage}
  \centering
  {\Large Universidad de Antioquia\par}
  \vspace{2cm}
  {\huge\bfseries Física Matemática I\par}
  \vspace{0.4cm}
  {\Large Notas de Clase Parte 1\par}
  \vspace{2.5cm}
  {\large Autor:\par}
  \vspace{0.2cm}
  {\large Daniel Soto Villada\par}
  {\large Estudiante de Astronomía\par}
  \vfill
  {\large \today\par}
\end{titlepage}

\newpage
\tableofcontents
\newpage

\section{Teoría de Transformación}

\subsection{Fundamentos de la Transformación de Coordenadas}
En un espacio euclidiano, es posible generar múltiples sistemas coordenados a partir de un sistema cartesiano base $(x, y, z)$ mediante la introducción de superficies definidas por $u_i = f(x, y, z)$. Recíprocamente, las coordenadas cartesianas se pueden expresar como funciones de las coordenadas curvilíneas: $x = x(u_i)$, $y = y(u_i)$ y $z = z(u_i)$.

A partir de estas relaciones, el diferencial del vector posición $\mathbf{r}$ se define mediante la regla de la cadena como:
\begin{equation}
d\mathbf{r} = \sum_{i=1}^{3} \frac{\partial \mathbf{r}}{\partial u_i} du_i
\end{equation}
donde el vector $\frac{\partial \mathbf{r}}{\partial u_i}$ es tangente a la curva coordenada $u_i$.

\subsection{Factores de Escala y Vectores Unitarios}
Para establecer una base normalizada de vectores unitarios $\hat{e}_i$, se introducen los factores de escala $h_i$, definidos por la magnitud del vector tangente:
\begin{equation}
h_i = \left| \frac{\partial \mathbf{r}}{\partial u_i} \right| \quad \Rightarrow \quad \hat{e}_i = \frac{1}{h_i} \frac{\partial \mathbf{r}}{\partial u_i}
\end{equation}
De este modo, el diferencial de línea se expresa como $d\mathbf{r} = \sum h_i \hat{e}_i du_i$. El texto asume el uso de bases ortonormales, donde el producto escalar entre vectores unitarios sigue la propiedad de la delta de Kronecker: $\hat{e}_i \cdot \hat{e}_j = \delta_{ij}$.

\subsection{Ejemplo: Transformación de Cartesianas a Esféricas}
Como aplicación práctica, se analiza el sistema de coordenadas esféricas $(r, \theta, \phi)$. Las relaciones de transformación son:
\begin{itemize}
    \item $x = r \sin\theta \cos\phi$
    \item $y = r \sin\theta \sin\phi$
    \item $z = r \cos\theta$
\end{itemize}

Calculando las derivadas parciales del vector posición respecto a cada coordenada y obteniendo sus magnitudes, se determinan los factores de escala esféricos:
\begin{equation}
h_r = 1, \quad h_\theta = r, \quad h_\phi = r \sin\theta
\end{equation}

Finalmente, los vectores unitarios de la base esférica se expresan en términos de la base cartesiana $(\hat{i}, \hat{j}, \hat{k})$ como:
\begin{align}
    \hat{e}_r &= \hat{i} \sin\theta \cos\phi + \hat{j} \sin\theta \sin\phi + \hat{k} \cos\theta \\
    \hat{e}_\theta &= \hat{i} \cos\theta \cos\phi + \hat{j} \cos\theta \sin\phi - \hat{k} \sin\theta \\
    \hat{e}_\phi &= -\hat{i} \sin\phi + \hat{j} \cos\phi
\end{align}
Esta transformación puede representarse mediante una matriz de rotación ortogonal $A$ cuyo determinante es igual a 1.


\section{El Símbolo de Levi-Civita y el Producto Vectorial}

\subsection{Definición Matemática (Ecuación 1.23)}
La ecuación (1.23) establece la regla para el producto vectorial entre dos vectores unitarios de una base ortonormal en un espacio euclidiano tridimensional:
\begin{equation}
    \hat{e}_i \times \hat{e}_j = \sum_{k=1}^{3} \epsilon_{ijk} \hat{e}_k
\end{equation}
Donde $\epsilon_{ijk}$ es el denominado **símbolo de Levi-Civita** o símbolo totalmente antisimétrico.

\subsection{Propiedades del Símbolo $\epsilon_{ijk}$}
De acuerdo con las fuentes, este símbolo se define bajo las siguientes reglas:
\begin{itemize}
    \item \textbf{Valor fundamental:} $\epsilon_{123} = 1$.
    \item \textbf{Antisimetría total:} El símbolo cambia de signo ante el intercambio de cualquier pareja de índices contiguos:
    \begin{equation}
        \epsilon_{ijk} = -\epsilon_{ikj} = \epsilon_{jki} = -\epsilon_{jik} = \epsilon_{kij} = -\epsilon_{kji}
    \end{equation}
    \item \textbf{Condición de anulación:} Si cualquier par de índices se repite (por ejemplo, $\epsilon_{112}$), el valor del símbolo es $0$.
\end{itemize}

\subsection{Contexto y Utilidad}
Esta expresión permite sintetizar de manera compacta las relaciones de producto cruz en sistemas de coordenadas curvilíneas ortogonales, tales como:
\begin{align*}
    \hat{e}_1 \times \hat{e}_2 &= \hat{e}_3 \\
    \hat{e}_2 \times \hat{e}_3 &= \hat{e}_1 \\
    \hat{e}_3 \times \hat{e}_1 &= \hat{e}_2
\end{align*}

\section{Operadores Diferenciales en Coordenadas Curvilíneas}

En un sistema de coordenadas curvilíneas ortogonales $(u_1, u_2, u_3)$ con factores de escala $h_i$ y vectores unitarios $\hat{e}_i$, las operaciones diferenciales básicas se definen de forma invariante [1, 2].

\subsection{Gradiente ($\nabla \phi$)}
El gradiente de un campo escalar $\phi(u_1, u_2, u_3)$ representa la tasa de cambio de la función en el espacio. Matemáticamente se define como:
\begin{equation}
    \nabla\phi \equiv \sum_{i=1}^3 \frac{\hat{e}_i}{h_i} \frac{\partial \phi}{\partial u_i}
\end{equation}
Propiedades clave:
\begin{itemize}
    \item El vector $\nabla\phi$ es perpendicular a las superficies equipotenciales ($\phi = cte$) [3].
    \item Su módulo corresponde al valor máximo de la derivada direccional, apuntando en la dirección de máximo crecimiento de la función [3].
\end{itemize}

\subsection{Divergencia ($\nabla \cdot \mathbf{B}$)}
La divergencia de un campo vectorial $\mathbf{B}$ se define como el flujo neto por unidad de volumen en un punto [4]. Su expresión general, utilizando $h = h_1 h_2 h_3$, es:
\begin{equation}
    \nabla \cdot \mathbf{B} = \frac{1}{h} \sum_{i=1}^3 \frac{\partial}{\partial u_i} \left( \frac{B_i h}{h_i} \right)
\end{equation}
Un campo donde $\nabla \cdot \mathbf{B} = 0$ se denomina \textbf{solenoidal} [5].

\subsection{Rotacional ($\nabla \times \mathbf{B}$)}
El rotacional (o \textit{curl}) mide la circulación máxima por unidad de área [6]. Utilizando el símbolo de Levi-Civita $\epsilon_{ijk}$, se expresa como:
\begin{equation}
    \nabla \times \mathbf{B} = \frac{1}{h} \sum_{i,j,k=1}^3 \hat{e}_i \epsilon_{ijk} h_i \frac{\partial}{\partial u_j} (B_k h_k)
\end{equation}
Un campo vectorial cuyo rotacional es nulo se define como \textbf{irrotacional} [7].

\subsection{Laplaciano ($\nabla^2 f$)}
El laplaciano de una función escalar es la divergencia de su gradiente ($\nabla \cdot \nabla f$) [8]. En coordenadas curvilíneas toma la forma:
\begin{equation}
    \nabla^2 f = \frac{1}{h} \sum_{i=1}^3 \frac{\partial}{\partial u_i} \left( \frac{h}{h_i^2} \frac{\partial f}{\partial u_i} \right)
\end{equation}
Este operador es fundamental en áreas como la electrostática, difusión de calor y mecánica cuántica [9].

\subsection{Identidades Fundamentales}
De las definiciones anteriores se derivan dos identidades críticas en la teoría de campos:
\begin{enumerate}
    \item \textbf{La divergencia de cualquier rotacional es nula:} $\nabla \cdot (\nabla \times \mathbf{A}) = 0$ [10].
    \item \textbf{El rotacional de cualquier gradiente es nulo:} $\nabla \times \nabla \phi = 0$ [7].
\end{enumerate}

\bibliographystyle{plainurl}
\bibliography{bibliografia} 
\end{document}
